% !TeX root = ../main.tex

\chapter{Analisi Economica: SLO, Costi e ROA}
\label{cap:analisi_economica}

\section{SLO e Costo del Downtime}
\label{sec:slo_downtime}

Sono stati definiti tre SLO coerenti col workload:
\begin{itemize}
    \item \textbf{SLO 1 -- Latenza assegnazione:} 90\% degli ordini \textit{high priority} assegnati entro 1 minuto (giustificato dal \texttt{time.sleep(30)} + latenza inferenza).
    \item \textbf{SLO 2 -- Disponibilità data layer:} 99{,}9\% di uptime per MQTT e InfluxDB; una perdita prolungata renderebbe cieco l'\texttt{HealthAgent}.
    \item \textbf{SLO 3 -- Affidabilità agentica:} 99\% di tool call valide senza innesco di \texttt{MAX\_AGENT\_ITERATIONS}.
\end{itemize}

Nel dominio della logistica critica, un'ora di downtime è stimata a circa \textbf{15.000\,\texteuro/ora} (media dei guadagni orari + penali contrattuali per SLA violati).

\section{Return on Agent (ROA): OPEX e CAPEX}
\label{sec:roa}

\textbf{Ottimizzazione OPEX.} (i) \textit{Triage preventivo}: il \texttt{triage\_manager} attiva gli agenti solo in caso di reale sbilanciamento, azzerando i cicli no-op. (ii) \textit{Prompt engineering}: i prompt vietano testo libero e impongono output esclusivamente tramite \textit{tool calling}, minimizzando i token di completamento. (iii) \textit{Monitoraggio consumi}: auditing sistematico dei token spesi per ogni ciclo.

\textbf{Preservazione CAPEX.} (i) \textit{Manutenzione preventiva dinamica}: il \texttt{LogisticAgent} include nei prompt la regola di evitare carichi $>3$\,kg su droni con usura $>70\%$ (\textit{soft constraint} prompt-driven). (ii) \textit{Elasticità cloud}: l'\texttt{HealthAgent} modula le repliche secondo il carico effettivo, riducendo l'over-provisioning. (iii) \textit{Operational Safety Cap}: limite di 6 droni in autonomia, oltre il quale interviene HITL. (iv) \textit{Throughput remunerativo}: la coda è ordinata per priorità, con il valore economico come criterio secondario.

\subsection*{Proiezioni Economiche}
Ipotizzando 2.880 cicli/giorno (polling ogni 30\,s):

\begin{table}[ht]
\centering
\begin{tabular}{lrr}
\toprule
\textbf{Voce} & \textbf{Gemini 2.5 Pro} & \textbf{Mistral Small} \\
\midrule
Infrastruttura cloud (1 control-plane + 2 worker) & 150{,}00 \$/mese & 150{,}00 \$/mese \\
Token LLM & $\sim$500{,}00 \$/mese & $\sim$10{,}00 \$/mese \\
\midrule
\textbf{OPEX totale mensile} & \textbf{$\sim$650{,}00 \$/mese} & \textbf{$\sim$160{,}00 \$/mese} \\
\bottomrule
\end{tabular}
\caption{Confronto OPEX al variare del modello LLM.}
\label{tab:costi_opex}
\end{table}

La forbice di costo (delta $\sim$490\,\$/mese) consente di modulare la scelta del modello: \textit{Mistral Small} per scenari ordinari di instradamento; \textit{Gemini 2.5 Pro} riservato a finestre critiche con anomalie complesse o conflitti di policy.

\section{Trade-off Costo/Affidabilità e Automazione/Sicurezza}
\label{sec:analisi_tradeoff}

Rispetto a un sistema deterministico (OPEX minimi ma fragile a eccezioni e dipendente dall'intervento manuale), il sistema multi-agente paga un costo variabile in token in cambio di flessibilità semantica: contestualizza gli errori e devia su flussi alternativi senza intervento umano. Sul piano automazione/sicurezza, il deterministico garantisce l'osservanza rigorosa delle policy ma è rigido davanti ad anomalie inattese; il multi-agente massimizza l'efficienza decisionale ma introduce rischi di allucinazione nelle tool call, mitigati dai guardrail del Capitolo~\ref{cap:safety_guardrails}.
